% =====================================================================
% บทที่ 5 การใช้งานสำหรับอาจารย์ที่ปรึกษาและผู้ดูแลรายวิชา
% =====================================================================
\label{ch:advisor}

\begin{guide}
บทของบทบาทถัดไปใช้รูปแบบเดียวกับบทที่ 4 งานที่ทำได้เหมือนกันทุกบทบาท
(เช่น การให้ความเห็น) ให้อธิบายไว้ที่เดียวแล้วอ้างถึงด้วยหมายเลขหัวข้อ
ไม่ต้องเขียนซ้ำในทุกบท
\end{guide}

\section{การพิจารณาข้อเสนอโครงงาน}

\begin{steps}
  \item ที่หน้ารายการโครงงาน เลือกแท็บ \btn{รอพิจารณา}
  \item คลิกข้อเสนอที่ต้องการเพื่อเปิดรายละเอียด
  \item พิมพ์ความเห็นในช่อง \field{ความเห็น}
  \item คลิก \btn{อนุมัติ} หรือ \btn{ส่งกลับแก้ไข}
        \result{ถ้าอนุมัติ ระบบสร้างโครงงานพร้อมไฟล์แม่แบบ 12 ไฟล์ให้นักศึกษา
        ถ้าส่งกลับ นักศึกษาจะได้รับการแจ้งเตือนให้แก้ไข}
\end{steps}

\begin{tip}
เขียนความเห็นให้ระบุสิ่งที่ต้องแก้เป็นข้อ ๆ นักศึกษาจะเห็นความเห็นนี้ที่หน้าข้อเสนอเมื่อถูกส่งกลับ
\end{tip}

\section{การให้ความเห็นในเอกสาร}
\label{sec:comment}

อาจารย์ที่ปรึกษาให้ความเห็นในไฟล์เอกสารของนักศึกษาได้โดยตรงผ่านแผงความเห็นดังภาพที่~\ref{fig:comment}
นักศึกษาตอบกลับได้ในแผงเดียวกัน

\begin{screenshot}{images/editor-comment.png}{แผงความเห็นใน LaTeX Editor}{fig:comment}
  \calloutline{0.70}{0.13}{0.60}{0.13}{1}
  \callout{0.665}{0.33}{2}
  \callout{0.665}{0.41}{3}
  \callout{0.665}{0.55}{4}
  \calloutline{0.862}{0.17}{0.862}{0.05}{5}
\end{screenshot}

\begin{calloutlist}
  \item ส่วนหัวแผงความเห็น แสดงจำนวนความเห็นและไฟล์ที่เปิดอยู่ กรองได้ตามสถานะ ทั้งหมด เปิด และแก้ไขแล้ว
  \item ช่องพิมพ์ความเห็นใหม่ เลือกข้อความใน Editor ก่อนเพื่อผูกความเห็นกับตำแหน่งนั้น
  \item ปุ่ม \btn{เพิ่มความเห็น} (หรือกด \keys{Ctrl+Enter})
  \item ความเห็นที่มีอยู่ พร้อมปุ่ม \btn{ตอบกลับ} และ \btn{แก้ไขแล้ว}
  \item ปุ่มเปิดหรือปิดแผงความเห็น ตัวเลขคือจำนวนความเห็นที่ยังเปิดอยู่
\end{calloutlist}

\begin{steps}
  \item เปิดโครงงานของนักศึกษา แล้วคลิกปุ่มหมายเลข~\calloutmark{5} เพื่อเปิดแผงความเห็น
  \item เลือกข้อความใน Editor ที่ต้องการให้ความเห็น
  \item พิมพ์ความเห็นในช่องหมายเลข~\calloutmark{2} แล้วคลิก \btn{เพิ่มความเห็น}
  \item เมื่อนักศึกษาแก้ไขตามความเห็นแล้ว คลิก \btn{แก้ไขแล้ว} เพื่อปิดความเห็น
\end{steps}

\section{การกำหนดวันส่งงาน}

อาจารย์ที่ปรึกษาและผู้ดูแลรายวิชากำหนดวันส่งงานของแต่ละรายการได้
นักศึกษาจะเห็นกำหนดส่งที่หน้าแรก และระบบจะแจ้งเตือนเมื่อใกล้ถึงกำหนด

% ตัวอย่าง: เขียนขั้นตอนตามหน้าจอจริงของระบบ พร้อมภาพหน้าจอ
% \begin{steps}
%   \item ...
% \end{steps}

\section{แดชบอร์ดสำหรับผู้ดูแลรายวิชา}
\label{sec:dashboard}

ผู้ดูแลรายวิชาและอาจารย์ที่ปรึกษาเปิดเมนูแดชบอร์ดได้ที่ \cmd{/dashboard} เพื่อดูจำนวนโครงงาน
ข้อเสนอที่รอพิจารณา และจำนวนการส่งงาน และส่งออกข้อมูลเป็นไฟล์ CSV ได้
อาจารย์ที่ต้องเป็นผู้ดูแลรายวิชาให้แจ้งผู้ดูแลระบบกำหนดบทบาทให้ (หัวข้อ~\ref{sec:users})
